% エントリ 8(2026-11-03): 天体物理の入力、X 線連星、(mu, f) 面の窓
\section{[E8, 2026-11-03] 窓を $(\mu,f)$ 面に描く: X 線連星と年齢の入力}\label{sec:e08}

E3--E7 で、窓は $f\sim10^{12}$--$10^{13}\GeV$ にありうるとわかった。ここでは、二つの外部の入力を重ねて窓を $(\mu,f)$ 面に描く。
\begin{itemize}
\item \textbf{X 線連星のスピン}: 同じ $(\mu,f)$ を独立に制限する。
\item \textbf{世代ごとの年齢}: 窓が開くかを決める。
\end{itemize}
抜き書きは \texttt{Refs/notes/window\_consistency.md}(ほかの観測との整合)と \texttt{Refs/notes/astro\_inputs.md}(年齢と 2 世代目のスピン)にある。計算は \texttt{calc/e08\_window.py}、出力は \texttt{calc/out/e08.json}。

\subsection{天体物理の入力(本文で照合)}

\paragraph{2 世代目の自然スピン}
\begin{itemize}
\item 終スピンの当てはめは \arxiv{1605.01938} の一般スピン版(式 12--17、Table 1)を使う。等質量・無スピンで 0.68646。誤差は整列スピンで rms 0.002、一般スピンで 0.006。
\item 終質量は \arxiv{1206.3803} の式 18 を使う。
\item 球状星団の模擬計算(\arxiv{1906.10260} の Fig.~8)では、1 世代目が回らないとき 2 世代目の 90\% が 0.67--0.69 に入る。1 世代目のスピンが 0.5 のときは中央値 0.62、90\% が 0.53--0.72 である。
\end{itemize}

\paragraph{世代ごとの年齢}
2 世代目の年齢(親の合体から自分の合体まで)の分布を直接示した論文はない。近い量は次のとおり(図からの読み取りを含む)。

\begin{table}[H]
\centering\small
\caption{世代ごとの年齢の目安(\texttt{Refs/notes/astro\_inputs.md})。}
\label{tab:e08_age}
\begin{tabularx}{\linewidth}{@{}l>{\raggedright\arraybackslash}X>{\raggedright\arraybackslash}X@{}}
\toprule
環境 & 2 世代目の年齢 & 1 世代目の年齢 \\
\midrule
球状星団 & 典型 10--100 Myr、下限 $\sim1$ Myr(\arxiv{2103.05016} Fig.~2 から読み取り、\arxiv{1811.03640} の沈降時間) & 低赤方偏移の合体で $\gtrsim6$ Gyr(\arxiv{1911.00018} Fig.~15) \\
核星団 & 典型 10--100 Myr、Gyr まで裾(\arxiv{2103.05016}、Fragione \& Silk 2020) & 記載なし \\
AGN 円盤 & $\lesssim0.01$--$0.1$ Myr(円盤の寿命 0.1--10 Myr のうち) & 記載なし \\
孤立連星 & -- & $p(\tau)\propto1/\tau$、GWTC-2 で中央値 2.8 Gyr \\
\bottomrule
\end{tabularx}
\end{table}

GW241011 について LVK(\arxiv{2510.26931})は、星団(球状・核・若い大質量)での 2 世代目の合体と解釈している。主星の年齢は推定していない。

\subsection{X 線連星の制限をこのノートの模型で計算する}

\arxiv{2412.03655} の Table III の値を使い、準平衡(E3)、相対論的な超放射の率(E4)、主要項の $\gamma_\mathrm{BH}$ で、次の条件を課した。最初のスピンは、測られた下限から最も遠い値に置いた(控えめな選び方)。
\begin{itemize}
\item \textbf{GRS 1915+105}: $12.4\Msun$、$\tilde a>0.95$、年齢 $\tau_\mathrm{Sal}\approx4.5\times10^7\yr$。スピンが 0.995 から 0.95 まで落ちない。
\item \textbf{Cyg X-1}(控えめなスピン): $21.2\Msun$、$\tilde a=0.92^{+0.05}_{-0.07}$、$4.8\times10^6\yr$。0.995 から 0.85 まで落ちない。
\item \textbf{Cyg X-1}(fiducial のスピン): $\tilde a>0.9985$。0.9995 から 0.9985 まで落ちない。$\chi>0.995$ は E4 の表の外なので、解析式の率を使った。
\end{itemize}

\paragraph{照合}
$\mu=0.69\times10^{-12}\eV$ での上限は次のとおり。図から読んだ \arxiv{2412.03655} の値と比べた(\texttt{window\_consistency.md})。
\begin{itemize}
\item GRS 1915+105: このノートで $f\le2.5\times10^{12}\GeV$、同論文の Fig.~1 は $2.3\times10^{12}\GeV$。同論文はベイズの 2$\sigma$ で、$n\le5$ の準位と相対論的な $\gamma_\mathrm{BH}$ を使っている。
\item Cyg X-1(控えめ): このノートで $5.3\times10^{12}$、同論文の $n\le3$ で $8.3\times10^{12}\GeV$。
\item Cyg X-1(fiducial): このノートで $2.9\times10^{11}$、同論文で $(0.75$--$0.78)\times10^{12}\GeV$。このノートは最初のスピンを 0.9995 に固定したので強めに出る。
\end{itemize}
GRS 1915+105 の値は 10\% で合う。

\subsection{結果}

\begin{figure}[H]
\centering
% E8: (mu, f) 面の窓。データは calc/e08_window.py(gamma_BH は主要項)
\begin{tikzpicture}
\begin{axis}[note, width=0.56\linewidth, height=0.5\linewidth, xmode=log, ymode=log,
  xmin=0.4e-12, xmax=1.2e-12, ymin=1e11, ymax=2e13,
  xtick={0.4e-12,0.5e-12,0.6e-12,0.7e-12,0.8e-12,1e-12,1.2e-12},
  xticklabels={0.4,0.5,0.6,0.7,0.8,1.0,1.2},
  xlabel={$\mu\ [10^{-12}\,\mathrm{eV}]$}, ylabel={$f\ [\mathrm{GeV}]$},
  legend style={font=\scriptsize, at={(1.02,1.0)}, anchor=north west, fill=white, draw=none}, legend cell align=left]
  % 2609.02678 の質量(90%)
  \fill[cB!10] (axis cs:0.60e-12,1e11) rectangle (axis cs:0.89e-12,2e13);
  % GRS 1915+105 の排除(上側)
  \addplot[name path=grs, ink, line width=1.0pt] table[x=mu, y=GRS] {figs/data/e08_window_lo.dat};
  \addlegendentry{GRS 1915+105(上を排除)}
  \addplot[name path=top, draw=none, forget plot] coordinates {(0.4e-12,2e13) (1.2e-12,2e13)};
  \addplot[cM!30, forget plot] fill between[of=grs and top];
  \addplot[ink, line width=0.8pt, dashed] table[x=mu, y=CygCons] {figs/data/e08_window_lo.dat};
  \addlegendentry{Cyg X-1、$\tilde a=0.92$(上を排除)}
  \addplot[ink, line width=0.8pt, densely dotted] table[x=mu, y=CygFid] {figs/data/e08_window_lo.dat};
  \addlegendentry{Cyg X-1、$\tilde a>0.9985$(上を排除)}
  % GW241011 が回転を保つ上限
  \addplot[cA!45, line width=0.9pt] table[x=mu, y=GW241011_T1e7] {figs/data/e08_window_lo.dat};
  \addlegendentry{GW241011 が保つ: $T_2=10^7\yr$}
  \addplot[cA, line width=0.9pt] table[x=mu, y=GW241011_T3e7] {figs/data/e08_window_lo.dat};
  \addlegendentry{\hphantom{GW241011 が保つ: }$3\times10^7\yr$}
  \addplot[cA!60!black, line width=0.9pt] table[x=mu, y=GW241011_T1e8] {figs/data/e08_window_lo.dat};
  \addlegendentry{\hphantom{GW241011 が保つ: }$10^8\yr$}
  % 1 世代目が 6 Gyr で回転を失う下限(chi0 = 0.7、delta = 0.18)
  \addplot[cC!60!black, line width=0.9pt, dashed] table[x=mu, y=M10_d18] {figs/data/e08_window_lo.dat};
  \addlegendentry{1 世代目が失う: $10\Msun$}
  \addplot[cC, line width=0.9pt, dashed] table[x=mu, y=M15_d18] {figs/data/e08_window_lo.dat};
  \addlegendentry{\hphantom{1 世代目が失う: }$15\Msun$}
  \addplot[cC!45, line width=0.9pt, dashed] table[x=mu, y=M20_d18] {figs/data/e08_window_lo.dat};
  \addlegendentry{\hphantom{1 世代目が失う: }$20\Msun$}
\end{axis}
\end{tikzpicture}

\caption{$(\mu,f)$ 面の窓($\gamma_\mathrm{BH}$ は主要項)。橙の帯は \arxiv{2609.02678} の質量(90\%)。灰色は GRS 1915+105 で排除。青の実線より下では GW241011 の主星が年齢 $T_2$ のあいだ 0.69 を保つ。緑の破線より上では、1 世代目($\chi_0=0.7$)が 6 Gyr のうちに 211 の臨界 $+0.18$ まで回転を失う。窓は、灰色と青の線の下で、かつ緑の線の上にある。}
\label{fig:e08}
\end{figure}

\begin{result}
$\mu=0.69\times10^{-12}\eV$ では次のようになる(主要項の $\gamma_\mathrm{BH}$。括弧は E7 の補正を入れた値)。
\begin{itemize}
\item \textbf{上限}: GRS 1915+105 で $f\le2.5\times10^{12}\GeV$($3.1\times10^{12}$)。GW241011 の主星は $T_2=10^7,3\times10^7,10^8\yr$ で $f\le4.4,2.5,1.4\times10^{12}\GeV$($7.0,4.0,2.2\times10^{12}$)。
\item \textbf{下限}: 1 世代目が 6 Gyr で回転を失うのに要る $f$ は、$\delta=0.18/0.05$ で次のとおり。
  \begin{itemize}
  \item $M_1=10\Msun$: $1.8/3.2\times10^{12}\GeV$
  \item $15\Msun$: $0.61/1.1\times10^{12}\GeV$
  \item $20\Msun$: $0.25/0.50\times10^{12}\GeV$
  \end{itemize}
\item \textbf{窓}: $f\approx(1$--$2.5)\times10^{12}\GeV$ の細い帯として残る。条件は三つある。
  \begin{enumerate}
  \item 2 世代目の年齢が $\lesssim3\times10^7\yr$ であること。球状星団の典型的な値(10--100 Myr)の下側にあたる。
  \item 積み上がりがおもに $M_1\gtrsim10$--$15\Msun$ に限られること。$10\Msun$ の 1 世代目は、$T_2\le3\times10^7\yr$ かつ $\delta=0.18$ のときだけ窓に入る。$\delta=0.05$ や $T_2=10^8\yr$ では窓の外になる。
  \item Cyg X-1 の fiducial のスピン($>0.9985$)が正しくないこと。正しければ $f\le2.9\times10^{11}\GeV$ になり、窓は閉じる。
  \end{enumerate}
\item \textbf{ほかの観測}(\texttt{window\_consistency.md}): GW231123、連続波・背景重力波の探索、銀河核での合体率は、この $f$ の範囲に届かない。超放射によらない制限は、標準模型との結合を仮定しない限りない。この $(\mu,f)$ の ALP は、ずれ機構では暗黒物質のすべてにはならない($f\gtrsim2\times10^{13}\GeV$ が要る)。
\end{itemize}
\end{result}

\paragraph{意味}
問いは三つの観測の重なりに絞られた。三つとも $f\approx2\times10^{12}\GeV$ の近くで効いている。
\begin{itemize}
\item 1 世代目の積み上がり(GWTC-5)。
\item 若い 2 世代目の速い回転(GW241011)。
\item 長く生きた X 線連星の速い回転(GRS 1915+105)。
\end{itemize}
集団解析では、2 世代目の年齢の分布と積み上がりの質量の範囲を同時に推定する。その結果が、窓を閉じるか、$f$ を特定するかを決める。Cyg X-1 のスピンの系統誤差は、外部の入力として明示する。

\begin{caveat}
\begin{itemize}
\item X 線連星の条件は決定論的な簡略版で、\arxiv{2412.03655} のベイズの扱い(最初のスピンと質量・年齢の周辺化)ではない。GRS の値が 10\% で合ったのは、手法の違いが小さいからとは限らない。
\item 2 世代目の年齢は、分布ではなく目安(図の読み取りを含む)である。
\item 1 世代目の判定($\chi_0=0.7$、$\delta$、$T_1=6$ Gyr)は説明のための選び方である。孤立連星の年齢(中央値 2.8 Gyr)を使うと、下限の $f$ は $\sqrt{6/2.8}\approx1.5$ 倍上がる。
\end{itemize}
\end{caveat}
